\chapter{Hand, Foot and Mouth Disease}

\section*{Definition and Overview}
Hand, foot and mouth disease (HFMD) is a common, mild viral illness that mainly affects young children, causing a fever, sores in the mouth, and a rash of small blisters on the hands, feet, and sometimes the buttocks. It is caused by a group of viruses called enteroviruses, most commonly coxsackievirus A16. The infection spreads easily through contact with saliva, fluid from blisters, faeces, and contaminated surfaces, and it is common in childcare and preschool settings. HFMD is not the same as foot-and-mouth disease, which affects animals. There is no specific treatment or vaccine, but the illness resolves on its own within 7--10 days, and care focuses on fluids, pain relief, and comfort.

\section*{Epidemiology}
Hand, foot and mouth disease is very common worldwide, mainly affecting children under five, though older children and adults can also be infected. It occurs throughout the year in Australia, with peaks in summer and autumn, and outbreaks are frequent in childcare centres. Most children have at least one episode. The coxsackievirus A6 strain has caused larger outbreaks with more extensive rashes in recent years. A more severe form, caused by enterovirus 71, is seen in parts of Asia but is uncommon in Australia.

\section*{Aetiology and Pathophysiology}
HFMD is caused by enteroviruses that infect the throat and intestine.

\begin{itemize}
    \item \textbf{Viruses:} most commonly coxsackievirus A16 and enterovirus 71, which belong to the enterovirus family
    \item \textbf{Transmission:} through direct contact with saliva, nasal mucus, fluid from blisters, and faeces, and via contaminated hands and surfaces
    \item \textbf{Incubation:} symptoms appear 3--7 days after exposure
    \item \textbf{Viral spread:} the virus enters through the mouth, multiplies in the throat and gut, and produces the characteristic rash and mouth ulcers
    \item \textbf{Infectious period:} people are most infectious in the first week, and the virus can be shed in faeces for weeks
    \item \textbf{Immunity:} infection produces immunity to that specific virus, but other strains can cause repeat infections
\end{itemize}

\section*{Clinical Features}
Symptoms usually begin with a fever and progress over a few days.

\begin{itemize}
    \item Fever, often the first sign
    \item Sore throat and mouth sores, which can make eating and drinking painful
    \item Small red spots that develop into blisters on the palms, soles, and sometimes the buttocks and knees
    \item Blisters that may be painful or itchy
    \item Reduced appetite, irritability, and generally feeling unwell
    \item In some children, a rash without mouth sores, or mouth sores without a rash
    \item In older children and adults, symptoms can be more severe
    \item Usually resolves within 7--10 days
\end{itemize}

\section*{Differential Diagnosis}
Other childhood illnesses with rashes and mouth sores should be considered.

\begin{itemize}
    \item Herpangina, a related enterovirus infection causing mouth ulcers without the hand and foot rash
    \item Chickenpox, which causes an itchy rash over the whole body
    \item Measles, which causes a blotchy rash and cough, coryza, and conjunctivitis
    \item Gingivostomatitis from the herpes simplex virus, which causes mouth ulcers and fever
    \item Impetigo and other skin infections
    \item Allergic reactions and other rashes
    \item Kawasaki disease, which is rare but serious
\end{itemize}

\section*{When to Seek Medical Care}
Most children with HFMD recover at home without medical care. See a doctor if the child is not drinking enough, shows signs of dehydration, has a high fever that persists, or is very unwell.

\textbf{Call 000 (or local emergency services) for:}
\begin{itemize}
    \item Signs of dehydration: little or no urine, dry mouth, sunken eyes, or lethargy
    \item A child who cannot keep down fluids
    \item A child who is unusually drowsy, confused, or difficult to wake
    \item Seizures or fits
    \item Signs of severe illness, including fast or laboured breathing
\end{itemize}

\section*{Diagnosis and Investigations}
HFMD is usually diagnosed clinically from the characteristic features.

\begin{itemize}
    \item \textbf{Clinical diagnosis:} the combination of fever, mouth sores, and the typical rash is usually sufficient
    \item \textbf{Throat, stool, or blister swabs:} PCR testing can confirm the virus in outbreaks or atypical cases
    \item \textbf{Blood tests:} rarely needed
    \item \textbf{Assessment of hydration:} checking for dehydration in affected children
    \item \textbf{Exclusion from childcare:} children stay home while unwell
\end{itemize}

\section*{Management}
Treatment is supportive, as the illness is viral and self-limiting.

\begin{itemize}
    \item \textbf{Fluids:} encourage frequent small drinks; cool fluids are often better tolerated than warm ones
    \item \textbf{Pain and fever relief:} paracetamol or ibuprofen in appropriate doses
    \item \textbf{Soft foods:} cool, bland foods such as yoghurt, ice cream, and jelly are easier to eat
    \item \textbf{Oral comfort:} avoid salty, acidic, or spicy foods that sting the mouth sores
    \item \textbf{Hydration monitoring:} watch for signs of dehydration, especially in young children
    \item \textbf{No antibiotics:} antibiotics do not work against viruses
    \item \textbf{Comfort measures:} rest and plenty of fluids
    \item \textbf{Review:} medical review if the child worsens or does not improve
\end{itemize}

\section*{Complications}
\begin{itemize}
    \item Dehydration from painful mouth sores limiting fluid intake
    \item Rarely, fingernail or toenail loss weeks after the illness, which grows back
    \item In rare cases, infection of the brain, heart, or lungs, mainly with enterovirus 71
    \item Spread to siblings, parents, and childcare contacts
    \item Pregnant women near term who contract HFMD should seek obstetric advice, though risk to the baby is low
\end{itemize}

\section*{Prognosis and Outcomes}
Hand, foot and mouth disease has an excellent prognosis. Symptoms settle within 7--10 days, and complications are rare. The mouth sores and rash heal without scarring, and the child returns to normal health. Immunity develops to the specific virus causing the infection, though other strains can cause repeat episodes. For the vast majority of children, HFMD is a mild, self-limiting illness that requires only comfort care and fluids.

\section*{Prevention and Self-Care}
\begin{itemize}
    \item Wash hands thoroughly with soap and water after nappy changes, toilet use, and contact with an infected child
    \item Clean and disinfect toys, surfaces, and utensils that may be contaminated
    \item Keep children with HFMD home from childcare and school while they are unwell
    \item Avoid close contact such as kissing and sharing cups with infected people
    \item Teach children good hand hygiene
    \item Do not pop blisters, and keep the skin clean
    \item Seek medical care if a child cannot drink enough or shows signs of dehydration
    \item Pregnant women should practise careful hygiene and seek advice if exposed
\end{itemize}

\section*{Sources and Further Reading}
The following reputable sources provide further information for healthcare students and patients seeking reliable, current advice about hand, foot and mouth disease.

\begin{itemize}
    \item Healthdirect Australia. \textit{Hand, foot and mouth disease.} https://www.healthdirect.gov.au/
    \item Australian Government Department of Health and Aged Care. \textit{Hand, foot and mouth disease fact sheet.} https://www.health.gov.au/
    \item NSW Health. \textit{Hand, foot and mouth disease fact sheet.} https://www.health.nsw.gov.au/
    \item Raising Children Network (Australian Government). \textit{Hand, foot and mouth disease.} https://raisingchildren.net.au/
    \item Centers for Disease Control and Prevention. \textit{Hand, foot, and mouth disease.} https://www.cdc.gov/
    \item NHS. \textit{Hand, foot and mouth disease.} https://www.nhs.uk/
\end{itemize}
